% Archived on 2026-07-31 when the active paper was shortened to retain only
% the adaptive Deterministic Schnorr result. This file is intentionally not
% input by main.tex and preserves the former static theorem as a tax file.
\section{Static-Message Unforgeability of the Derandomized Scheme} \label{app-der-static}

This appendix states and proves the polynomial-assumption result announced in Section~\ref{sec-derand}. The scheme is that of Construction~\ref{constr-der}, with or without the pre-hash; we write it without, since the pre-hash plays no role when no guessing is needed.


\begin{theorem} \label{thm-der-static}
Let $\GG$ be a group of prime order $p$ with generator $g$, and let $f\colon\GG\to\Zp$ be efficient with $\mathsf{Size}_f(0)=0$. Let $\advA=(\advA_0,\advA_1)$ be an adversary that runs in time $t_{\advA}$ and declares $q_s$ signing messages, where $q_s$ is not bounded by any parameter of the scheme. There are an adversary $\advB_{\mathsf{main}}$ against $\CDL[\GG,g,f]$, an adversary $\advB_{\mathsf{trig}}$ against DL, a distinguisher $\advD_0$ against $\PRF_0$, and, for $j\in[q_s]$ and $\beta\in\{0,1\}$, distinguishers $\advD^{\mathsf{io}}_{j,\beta}$ against $\iO$ and $\advD^{\mathsf{prf}}_{j,\beta,i}$ against $\PRF_i$ for $i\in[4]$, together with $\advD^{\mathsf{io}}_{\mathsf{trig}}$ and $\advD^{\mathsf{prf}}_{\mathsf{trig}}$, each running in time $t_{\advA}+\poly(\secp,q_s)$, such that, writing
$$
\epsilon_j
:=
\sum_{\beta\in\{0,1\}}
\Big(
\AdvIO{}{\advD^{\mathsf{io}}_{j,\beta}}
+\sum_{i=1}^{4}\AdvPRF{\PRF_i}{\advD^{\mathsf{prf}}_{j,\beta,i}}
\Big),
$$
we have
$$
\begin{aligned}
\AdvEUFSTATIC{\SchDScheme[\GG,\HF^{\mathsf{der}}]}{\advA}
\le\ &\AdvPRF{\PRF_0}{\advD_0}
+\AdvCDL{\GG,g,f}{\advB_{\mathsf{main}}} \\
&+\AdvDL{\GG,g}{\advB_{\mathsf{trig}}}
+\AdvIO{}{\advD^{\mathsf{io}}_{\mathsf{trig}}}
+\AdvPRF{\PRF_4}{\advD^{\mathsf{prf}}_{\mathsf{trig}}} \\
&+\sum_{j=1}^{q_s}\epsilon_j
+\frac{2q_s}{p} \;.
\end{aligned}
$$
\end{theorem}

\begin{corollary} \label{cor-der-static}
If $\CDL[\GG,g,f]$ holds, $\PRF_0,\ldots,\PRF_4$ and $\iO$ are secure, and $p$ is super-polynomial, then $\SchDScheme[\GG,\HF^{\mathsf{der}}]$ is EUF-static secure against every PPT adversary.
\end{corollary}

\paragraph*{Discussion.} Corollary~\ref{cor-der-static} is ordinary static security of one fixed scheme against every PPT adversary, with no query parameter anywhere and with every assumption polynomially hard. Compared with Corollary~\ref{cor-uf} it removes the bound $Q$. Compared with Corollary~\ref{cor-der-adaptive}, it uses only polynomially hard assumptions but does not make the queries adaptive.

Two features of the static game do the work. The declared messages are fixed before key generation, so the reduction knows every slot it will need to program and guesses nothing; this is what keeps the loss polynomial. And the trigger branch recomputes programmed points instead of storing them, so no table appears in the published circuit; this is what removes $Q$.

\begin{proof}[of Theorem~\ref{thm-der-static}]
Let $\advA_0$ declare a list of $q_s$ signing messages and a forgery message $m^*$ outside that list. Let $N\le q_s$ be the number of distinct signing messages, and relabel them $m_1,\ldots,m_N$ in order of first occurrence. Signing is deterministic, so every repeated position receives the signature of its first occurrence; the hybrid is indexed by these $N$ distinct messages.

\heading{Game $\Gm_0$ (uniform nonces).} The game replaces each $r_i=\PRFEv_0(k_0,m_i)$ by an independent uniform $r_i\getsr\Zp$. The key $k_0$ is used nowhere else, and the messages are declared before key generation, so a single distinguisher $\advD_0$ bounds the change by $\AdvPRF{\PRF_0}{\advD_0}$. From here on the honest nonces are uniform, as in the randomized scheme.

\heading{Games $\Gm_1,\ldots,\Gm_N$ (trigger signing).} For $0\le i\le N$, let $\Gm_i$ answer the first $i$ distinct declared messages by the trigger procedure, returning $(W_{m_j},s_{m_j})$, and the remaining distinct messages honestly; repeated positions replay the response of their first occurrence. Every game publishes the honest circuit $\CircDer$ of Figure~\ref{fig-circuit-derand}; hardwired circuits occur only inside the reductions between adjacent games.

Fix $j$ and compare $\Gm_{j-1}$ with $\Gm_j$. Both endpoints first introduce the event $\mathsf{Bad}_{\mathsf{eq}} := [R_j = W_{m_j}]$ and output $0$ when it occurs, where $R_j=g^{r_j}$ is the honest nonce point. Conditioned on the view before the change, $R_j$ is uniform and independent of $W_{m_j}$, so the two aborts cost at most $2/p$.

The rest of the step is the direct specialization of the local chain in Appendix~\ref{app-derand}. There is one simplification: the reduction is given $m_j$ by $\advA_0$ before key generation, so it does not guess. Concretely, it sets
$$
\begin{aligned}
t&=m_j,&
R&=g^{r_j},&
t_R&=\enc{R,m_j},\\
W&=g^{s_{m_j}}X^{-c_{m_j}},&
t_W&=\enc{W,m_j}.
\end{aligned}
$$
Let $v$ be the main-branch hash value at $(R,m_j)$. The reduction publishes the temporary two-point circuit from that chain, retyped so that its slot is the message. The circuit carries the two hash input--output pairs $\{(t_R,v),(t_W,c_{m_j})\}$, punctures $k_1,k_2$ at $T=\{t_R,t_W\}$, and punctures $k_3,k_4$ at $t$. It then runs the same exchange: make the $\PRF_1$ and $\PRF_2$ values on $T$ uniform, replace $v$ by a uniform value and discard the masks, change variables so the honest record is $(g^{s_0}X^{-v},v,s_0)$ for uniform $v,s_0$, make the $\PRF_3$ and $\PRF_4$ values at $t$ uniform, exchange the two records by Lemma~\ref{lem-swap}, and reverse. Each of the two directions invokes $\iO$ once and the four puncturable PRFs once, giving $\epsilon_j$. The slot $m_j$ and uniform nonce point $R$ are fixed before key generation. The trigger point $W$ depends on trigger-PRF outputs, but in each puncturable-PRF reduction the challenge inputs are sampled independently of the particular PRF key being challenged; the reduction samples the other independent keys after receiving those challenge values. Thus Definitions~\ref{def-pprf} and~\ref{def-io} apply with polynomial-time samplers. Hence
$$
\big|\Pr[\Gm_{j-1}\Rightarrow1]-\Pr[\Gm_j\Rightarrow1]\big|
\le \epsilon_j+\frac{2}{p} \;.
$$

In $\Gm_N$ every signature is a trigger signature, and computing $W_{m}$ and $s_{m}$ uses $X$ and the PRF keys but not $x$.

\heading{A main-branch forgery.} Let $E_{\mathsf{main}}$ be the event that $\advA_1$ returns a valid forgery on $m^*$ whose hash value comes from the main branch. Adversary $\advB_{\mathsf{main}}$ receives a CDL challenge $h$, sets $X\gets h$, samples the four PRF keys, publishes the obfuscated circuit, and answers the declared messages by the trigger procedure. Given a valid forgery $(R^*,s^*)$ it computes $a^*,b^*$ from $k_1,k_2$ at $\enc{R^*,m^*}$ and returns
$$
\big(R^*X^{a^*}g^{b^*},\ (s^*+b^*)\bmod p\big),
$$
which wins the CDL game by the computation of Section~\ref{sec-io-constr}, using that $f$ has no zeros. Hence $\Pr[E_{\mathsf{main}}]\le\AdvCDL{\GG,g,f}{\advB_{\mathsf{main}}}$.

\heading{A trigger-branch forgery.} Let $E_{\mathsf{trig}}$ be the event that the forgery is valid and its hash value comes from the trigger branch, so that $R^*=W_{m^*}$. Here the static setting removes the guessing used in Appendix~\ref{app-derand}: the forgery message $m^*$ is declared by $\advA_0$ before key generation, so the reduction knows the slot it must program.

Adversary $\advB_{\mathsf{trig}}$ receives a DL challenge $Y$, samples $x\getsr\Zp$ and sets $X\gets g^x$, and punctures $k_4$ at $m^*$. It publishes the trigger-forgery circuit used in Appendix~\ref{app-derand}, retyped so that the slot is the message. On input $(P,m^*)$, this circuit sets $c_{m^*}=\PRFEv_3(k_3,m^*)$, returns $c_{m^*}$ if $P=Z X^{-c_{m^*}}$, and otherwise runs the main branch of $\CircDer$; on slots different from $m^*$ it runs $\CircDer$ using the punctured key. With $Z=g^{\widehat s}$ for $\widehat s=\PRFEv_4(k_4,m^*)$ this circuit computes the same function as $\CircDer$, so $\iO$ justifies the switch. Puncturable-PRF security at $m^*$ replaces $\widehat s$ by a uniform value, after which $Z=Y$ has the right distribution. Signing uses the punctured key and never needs slot $m^*$, since $m^*$ is not among the declared messages. If the forgery takes the trigger branch at $m^*$ then
$$
g^{s^*}=\big(YX^{-c_{m^*}}\big)X^{c_{m^*}}=Y,
$$
so $\advB_{\mathsf{trig}}$ returns $s^*$ as the discrete logarithm of $Y$. Hence
$$
\Pr[E_{\mathsf{trig}}]\le
\AdvDL{\GG,g}{\advB_{\mathsf{trig}}}
+\AdvIO{}{\advD^{\mathsf{io}}_{\mathsf{trig}}}
+\AdvPRF{\PRF_4}{\advD^{\mathsf{prf}}_{\mathsf{trig}}} \;.
$$

\heading{Conclusion.} A valid forgery takes one of the two branches. Summing the $N$ adjacent-game bounds and using $N\le q_s$ gives the stated sum and the term $2q_s/p$; for notational consistency, the distinguishers for indices $j>N$ may be taken to be trivial. Every reduction runs in time $t_{\advA}+\poly(\secp,q_s)$, and no factor depending on the message space or on a guess appears.
\end{proof}
